\section{Benchmarks}
\label{sec:grust:benchmarks}

\GRust is designed to bridge the gap between high-level, model-based design and the safety,
performance, and concurrency features of \Rust that are critical for modern \sdv architectures.
It achieves this by unifying a deterministic, synchronous paradigm with high-level constructs for
asynchronous, event-driven behavior. A comprehensive evaluation must therefore validate the
effectiveness of this hybrid approach. Accordingly, our assessment is structured into three parts,
each targeting a core contribution. The evaluations were conducted during two internships that I
proposed, and their results provide independent validation for the contributions of this thesis.

The first two parts of the evaluation stem from an academic internship, conducted by Dimitry Doche,
and focused on validating the core principles of \GRust.
%
First, to assess the performance of the synchronous kernel, a quantitative analysis compares the
runtime efficiency of code generated by \GRust against that of established \Lustre compilers.
%
Second, to demonstrate improvements in design clarity, a qualitative case study contrasts a
high-level \GRust model of a complex concurrent system with its conventional low-level \Ci
implementation. Together, these results establish the performance and expressive power of the
language in a controlled academic setting.

The third and final part of our assessment demonstrates the practical applicability of \GRust for
\sdv development. This was the focus of an industrial internship at \ampere, conducted by Matthieu
Massif, which involved reimplementing existing industrial services in \GRust. The code metrics and
performance of these services were then compared on target hardware against the original,
handwritten asynchronous \Rust versions. A key outcome of this effort was the development of a
complete compilation and deployment toolchain, proving the integration of \GRust into a real-world
engineering workflow. These findings provide compelling evidence for the industrial viability of
\GRust, showing that it can be effectively integrated into the \Rust ecosystem at \ampere to
simplify the development of complex automotive software.

\subsection{\GRSync Performance Evaluation}
\label{sec:grust:benchmarks:compare_grsync}

%% Benchmark framework
This first evaluation assesses the performance of the synchronous kernel of \GRust,
\GRSync. To ensure a fair and direct comparison, we focus on purely synchronous systems,
which allows us to benchmark against established academic \Lustre compilers:
\Heptagon~\cite{heptagon}, \LustreC~\cite{lustrec}, and \Velus~\cite{velus}.
%% Goal: show that GSync is performant
The goal of this evaluation is to show that binaries generated by \GRSync for synchronous systems
relevant to our context are comparable behavior-wise and performance-wise to that of a
\Lustre-specific compiler. We propose to do so by comparing the run time of the binaries generated
by the compilers on a few realistic systems, using a series of 100,000 steps without interruption%
---meaning step $n+1$ starts as soon as step $n$ is over.

\subsubsection{Benchmark Tooling}
\label{sec:grust:benchmarks:compare_grsync:tools}

As part of the \GRust project, two key tools were developed by Adrien Champion in order to conduct
this evaluation: a \Lustre front-end for the \GRust compiler and a specialized benchmarking utility.
They are described in this section.

%% GLustre: Lustre => GRSync
\paragraph{\GLustre front-end}
To make \GRSync fit in this evaluation, we used a front-end for the \GRust compiler called
\GLustre that supports an expressive enough subset of the \Lustre language. It
front-end parses a \Lustre file, transpiles it to \GRSync, compiles the resulting file to \Rust code
and finally compiles the \Rust code. We thus obtain a binary for each of the \Lustre systems we want
to compare with.

%% Benchmark tool
\paragraph{\BenchIt tool}
To properly run the evaluation, we used \BenchIt, a tool specialized for
benchmarking synchronous systems by efficiently feeding them a large number of inputs and measuring
the time they take to produce the same number of step-outputs. Benchmarks are given as a \Lustre
file \texttt{<file>.lus} and an inputs file \texttt{<file>.toml} specifying a trace of inputs.
For a given benchmark, \BenchIt compares \Lustre compilers by: compiling the system with each
compiler to obtain a binary, pre-processing the inputs for each binary (they all have specific input
formats), measuring the time they take to produce all their step-outputs, and logging the run
times after confirming the traces of step-outputs.
Currently, \BenchIt supports the following compilers: \Heptagon, \GLustre, \LustreC, and \Velus.

Most \Lustre compilers produce a binary that starts with an initialization phase that dynamically
allocates all the memory it will need once it actually starts. We ignore the initialization phase
by waiting 500 milliseconds between starting the binaries and providing step-inputs, and
measure only the run time after we start feeding the inputs.

\BenchIt is also able to generate traces of random inputs for \Lustre systems. Ideally we would
want large, \emph{realistic} traces of inputs, but such traces are rarely available and
quite difficult to generate for embedded systems, usually requiring a complex, often
proprietary, simulation environment.

\subsubsection{Selected Systems}
\label{sec:grust:benchmarks:compare_grsync:selected_systems}

Given our industrial context, we are interested in reasonably realistic, ideally semi-industrial
systems. Also, we want the computation steps to be relatively expensive to lower the impact of the
low-level details specific to each compiler, such as the different ways the resulting binaries read
(write) their inputs (outputs).
%
For our evaluation, we selected several systems with increasing complexity that are presented below.
As \Lustre compilers have a tendency to differ in the \Lustre subset they support (state machines,
custom types, clocks, \etc), we had to rewrite and select the systems in the intersection of the
features and syntax accepted by all compilers to make sure all compilers work on the exact same
\Lustre source files.

\paragraph{Pololu}
%% Pololu %%
The \href{https://www.pololu.com/docs/0J83}{\pololu} is a small mobile platform designed for
educational purposes. It is controlled by an Arduino-type board, with reflectance sensors that
measure the grayscale intensity of the ground to the left, center, and right of the robot. Dimitry
Doche has implemented in \Lustre a simple controller that allows the robot to follow a dark line (by
comparing the measured grayscale to the expected grayscale of the line). As a sanity check, we have
embedded the \Ci code compiled from \Lustre: the controller step function is called in a \Ci
interface that connect the hardware to the controller. We computed the position of the robot
approximated from the speed of the wheels produced by the compilers on a randomly generated trace of
inputs. The results in \Cref{fig:grust:benchmarks:compare_grsync:selected_systems:pololu} confirm
that the \GLustre binary behaves consistently with the other three compilers, as shown by the
superposition of the four trajectories.

\begin{figure}[ht!]
    \centering
    \includegraphics[width=0.6\textwidth]{trajectory_comparison_better.png}
    \caption[Positions of the Pololu.]%
    {Positions of the Pololu computed from one input trace, generated by the \Lustre
        compilers (\Heptagon, \LustreC, \Velus, and \GLustre) in comparison with the position
        computed by the robot.}
    \label{fig:grust:benchmarks:compare_grsync:selected_systems:pololu}
\end{figure}

\paragraph{Crazyflie}
%% Crazyflie %%
The \href{https://www.bitcraze.io/products/crazyflie-2-1-plus/}{\crazyflie} is a quadcopter drone
shipping with an open source development platform. We are interested in its Kalman filter, a
statistics estimator that uses sensor measurements over time to compute an estimation of the current
state. The original implementation is provided as \Ci code in the
firmware~\cite{DBLP:conf/icra/MullerHD15,MuellerCovariance2016}.%
\footnote{
    The code is available at
    \url{https://github.com/bitcraze/crazyflie-firmware/blob/master/src/modules/src/estimator/estimator_kalman.c}.
}
Dimitry Doche, assisted by Basile Pesin and Timothy Bourke (who have reimplemented the PID
controller of the \crazyflie in \Velus), translated the core of the estimator into the
\Lustre dialect of \Heptagon, which supports arrays
and composite data types. The resulting program contains approximately 1,000 lines of code or
35,000 tokens, excluding comments. As a sanity check, we verified that the resulting estimator and
the initial firmware \Ci code was producing the same outputs using a real input trace. The trace
was reconstructed by concatenating data from multiple \crazyflie drone flights. Sensor measurements
were obtained \via radio transmission, but due to bandwidth limitations, it was not possible to log
the entire input stream at once. We then transpiled this program into a simpler \Lustre variant
compatible with all compilers, which involved inlining arrays and structures. This significantly
increased the number of tokens (175,000) and of lines of code (15,000). We verified that all four
compilers produce identical output traces for several input traces of 100,000 steps.

\paragraph{TCM}
%% TCM %%
The \emph{Transport Class Model} (TCM) is a control system for a generic, mid-size, twin-engine,
commercial transport-class aircraft. It was originally developed by \textsc{NASA} Langley in
\Simulink~\cite{tcmOriginal}. It was later translated to \Lustre and equipped with a full formal
specification based on the \emph{Federal Aviation Regulations}~\cite{tcmCocospec}.%
\footnote{
    Code available at
    \href{https://github.com/kind2-mc/cocospec_tcm_experiments/tree/3b0834704e8434f1da08e45f79f129330faca6fc/systems}
    {https://github.com/kind2-mc/cocospec\_tcm\_experiments}.
}
As a semi-industrial system, it is an ideal candidate to evaluate the performance of synchronous
systems compiled by \GRust.

\paragraph{Kind~2 benches}
%% kind2 benches %%
Next, we ran our evaluation on approximately 200 \Lustre programs from the benchmark suite of the
Kind~2 model-checker~\cite{DBLP:conf/cav/ChampionMST16} (those included in the subset of \Lustre
that we consider for our comparison). Even though several of these systems are rather small, leaving
not much room for improvement, \GLustre shows suspiciously good performance compared to other
\Lustre compilers. To investigate this, we added a control benchmark, \texttt{in-out.lus}, which
takes several inputs and simply outputs the same constant at each step. Since \GLustre remains
significantly faster on this benchmark, we infer that the performance difference is not only
explained by the way \GLustre compiles the systems, but rather most likely influenced by low-level
aspects of the \Ci and \Rust runtimes, such as input scanning and output printing.

\subsubsection{Evaluation and Discussion}

\begin{figure}[ht!]
    \centering
    \begin{subfigure}{1\textwidth}
        \centering
        \begin{longtable}[c]{|lcc | c | c | c | c |}
            \hline
                                       &       &     & in-out & Pololu & TCM    & Crazyflie \\\hline\hline
            \multirow{2}{*}{\GRSync}   & Mean  & (s) & 0.1445 & 0.126  & 0.2466 & 0.72      \\
                                       & $\pm$ & (s) & 0.0003 & 0.002  & 0.0005 & 0.03      \\\hline
            \multirow{2}{*}{\Heptagon} & Mean  & (s) & 0.336  & 0.233  & 0.5871 & 2.15      \\
                                       & $\pm$ & (s) & 0.002  & 0.003  & 0.0006 & 0.06      \\\hline
            \multirow{2}{*}{\LustreC}  & Mean  & (s) & 0.4454 & 0.247  & 0.615  & 2.54      \\
                                       & $\pm$ & (s) & 0.0004 & 0.004  & 0.002  & 0.07      \\\hline
            \multirow{2}{*}{\Velus}    & Mean  & (s) & 0.342  & 0.231  & 0.5890 & 2.16      \\
                                       & $\pm$ & (s) & 0.003  & 0.003  & 0.0003 & 0.06      \\\hline
        \end{longtable}
        \caption[Comparison of execution time.]%
        {Comparison of execution time (in seconds) for the \texttt{in-out.lus}, Pololu, TCM,
            and \crazyflie systems.}
        \label{table:grust:benchmarks:compare_grsync:selected_systems:exec-time:others}
    \end{subfigure}

    \begin{subfigure}[b]{1\textwidth}
        \centering
        \includegraphics[width=0.6\textwidth]{scatter_grust_comparaison_better.png}
        \caption[Comparison of execution time for the Kind~2 benches.]%
        {Comparison of execution time for the Kind~2 benches.}
        \label{fig:grust:benchmarks:compare_grsync:selected_systems:exec-time:kind2}
    \end{subfigure}
    \caption[Execution times for codes from \GRSync and academic compilers.]%
    {Combined comparisons of execution time on the different selected systems with input traces of
        100,000 steps, running each benchmark between 10 and 100 times.}
    \label{fig:grust:benchmarks:compare_grsync:selected_systems:exec-time}
\end{figure}

The performance evaluation, conducted by Dimitry Doche with assistance from Christophe Garion,
measured the execution time of the selected systems. The results are presented in
\Cref{fig:grust:benchmarks:compare_grsync:selected_systems:exec-time}. Each system was executed on
randomly generated input traces of 100,000 steps, and each benchmark was run between 10 and 100
times to ensure stable measurements.%
\footnote{
    The evaluation was performed on an Asus Zenbook 14 equipped with an AMD Ryzen 7 8840HS processor
    running Manjaro Linux.
}
We verified that all compilers produced equivalent output traces for each system, with the exception
of integer overflows, since \LustreC relies on 32-bit integers while the other compilers use 64-bit
integers.

\paragraph{Results}
The results for the evaluation on the bigger systems, detailed in
\Cref{table:grust:benchmarks:compare_grsync:selected_systems:exec-time:others}, show that \GRSync
(\via the \GLustre front-end) consistently outperforms the other compilers. For the most complex
systems, the TCM and the Kalman filter of the \crazyflie, the code generated by \GRSync runs two to
three times faster than the code from the next fastest compiler.

This trend is confirmed by the evaluation on the Kind~2 benchmark suite, illustrated in
\Cref{fig:grust:benchmarks:compare_grsync:selected_systems:exec-time:kind2}. The scatter plot
compares the execution time of \GRSync ($y$-axis) against the other compilers ($x$-axis) on a
logarithmic scale. To each point corresponds a benchmark program from Kind~2. A point below the main
diagonal ($y=x$) indicates that \GRSync is faster (the execution time of the \GRSync version
is smaller than the compared one). All the 200 benchmarks lie well below this diagonal,
demonstrating that the performance advantage \emph{of the toolchain of \GRSync} is consistent across
a wide variety of programs, not just a few large examples. This performance in suspiciously good,
these systems are rather small which leaves not much room for improvement.

The \texttt{in-out.lus} benchmark, which performs minimal computation, acts as a control to isolate
framework overhead. Even here, \GRSync is more than twice as fast as its counterparts
(\cf \Cref{table:grust:benchmarks:compare_grsync:selected_systems:exec-time:others}).

\paragraph{Analysis}
Even though the \GRSync toolchain is faster according to this evaluation, we \emph{do not claim}
that this is due to the way \GRSync compiles synchronous systems. As steps are often measured in
milliseconds, the low-level, \Ci/\Rust-specific details of scanning and parsing inputs, printing
outputs, operations on primitive types, \etc can become significant. This is highlighted by the
\texttt{in-out.lus} file, that operates as a scientific control; the evaluation of all the compilers
using this file showed a reduction of 200 milliseconds of the execution time for the code generated
by \GRSync.

However we \emph{do claim} that this evaluation shows that the synchronous part of \GRust and its
compilation framework relying on \Rust are competitive relative to state-of-the-art academic tools
that are specialized in `\Lustre to \Ci' compilation, despite \GRust targets a much richer language.

\subsection{Evaluation of the Design of a UAV with \GRust}
\label{sec:grust:benchmarks:eval_grust_design}

To demonstrate the benefits of using \GRust as a high-level modeling tool, we present a case study
centered on the Kalman filter from the state estimator of the \crazyflie UAV~\cite{crazyflie-fm},
described previously.

\paragraph{The original Kalman filter}
The original implements an adapted
\href{https://github.com/bitcraze/crazyflie-firmware/blob/b9d676aa3544b68eea17d11b0fbec437da909878/src/modules/src/kalman_core/kalman_core.c}
{Kalman filter}, which updates the estimated state of the drone based on incoming sensor
measurements~\cite{DBLP:conf/icra/MullerHD15,MuellerCovariance2016}.%
\footnote{
    Each important function of the estimator is referred with a link to the \Ci functions of the
    \crazyflie firmware.
}

The global algorithm of the stabilizer proceeds the following computations:
\begin{enumerate}
    \item \textbf{Reset:} The estimated state is reverted to its initial values upon system
          initialization or when a reset command is received.
    \item \textbf{Prediction:} The last estimated state is given to the physical model to
          \href{https://github.com/bitcraze/crazyflie-firmware/blob/b9d676aa3544b68eea17d11b0fbec437da909878/src/modules/src/kalman_core/kalman_core.c#L744C1-L791C2}
          {predict} the current state.
    \item \textbf{Transformation:} \href{https://github.com/bitcraze/crazyflie-firmware/blob/b9d676aa3544b68eea17d11b0fbec437da909878/src/modules/src/kalman_core/kalman_core.c#L613C1-L619C2}
          {Noise is added} to take into account noisy sensors, so that the estimated state is
          comparable to the measures.
    \item \textbf{Update:} If measurements are received, the parameters of the filter (covariance
          and innovation) are
          \href{https://github.com/bitcraze/crazyflie-firmware/blob/b9d676aa3544b68eea17d11b0fbec437da909878/src/modules/src/estimator/estimator_kalman.c#L299C1-L365C2}
          {updated} according to the estimated state and the measures.
    \item \textbf{Correction:} The new estimated state is
          \href{https://github.com/bitcraze/crazyflie-firmware/blob/b9d676aa3544b68eea17d11b0fbec437da909878/src/modules/src/kalman_core/kalman_core.c#L621C1-L742C2}
          {finalized} using the parameters of the filter and the error between the noisy estimated
          state and the measured state.
    \item \textbf{Externalization}: Finally, the internal state of the Kalman filter
          \href{https://github.com/bitcraze/crazyflie-firmware/blob/b9d676aa3544b68eea17d11b0fbec437da909878/src/modules/src/kalman_core/kalman_core.c#L744C1-L791C2}
          {is converted and copied} into the external state and format expected by other system
          modules.
\end{enumerate}

The firmware is built upon \textit{FreeRTOS}, an open-source real-time operating system that manages
concurrency through a predictable scheduler~\cite{freertos}. The estimation logic is not executed
directly within the main control loop of the stabilizer. Instead, it runs in a dedicated
\href{https://github.com/bitcraze/crazyflie-firmware/blob/b9d676aa3544b68eea17d11b0fbec437da909878/src/modules/src/estimator/estimator_kalman.c#L214}
{estimation task}, which is synchronized with the primary
\href{https://github.com/bitcraze/crazyflie-firmware/blob/b9d676aa3544b68eea17d11b0fbec437da909878/src/modules/src/stabilizer.c#L326}
{stabilizer task}.

The stabilizer task, which loops at 1,000~Hz, orchestrates the execution at two levels. At each
cycle, it first protects access to the shared state data by acquiring a semaphore, which preempts
the estimation task. After
\href{https://github.com/bitcraze/crazyflie-firmware/blob/b9d676aa3544b68eea17d11b0fbec437da909878/src/modules/src/estimator/estimator_kalman.c#L293}
{copying} the latest estimated state and freeing the data-protective semaphore, it releases a
\emph{different}
\href{https://github.com/bitcraze/crazyflie-firmware/blob/b9d676aa3544b68eea17d11b0fbec437da909878/src/modules/src/estimator/estimator_kalman.c#L296}
{semaphore}, signaling the estimation task to
\href{https://github.com/bitcraze/crazyflie-firmware/blob/b9d676aa3544b68eea17d11b0fbec437da909878/src/modules/src/estimator/estimator_kalman.c#L223}
{begin its next computation}. This intricate mechanism allows the Kalman filter to run concurrently
with the main
\href{https://github.com/bitcraze/crazyflie-firmware/blob/b9d676aa3544b68eea17d11b0fbec437da909878/src/modules/src/stabilizer.c#L349}
{controller} computations, ensuring the controller always uses the most recently available state.

However, this approach introduces significant complexity. The control flow is no longer sequential
but is dictated by the subtle interplay of tasks and semaphores. Understanding the system requires
reasoning about shared memory, potential race conditions, and the precise semantics of the
scheduler. Such low-level manual synchronization is not only difficult to reason about but is also a
common source of hard-to-debug concurrency errors, such as deadlocks or priority inversions. This
complexity motivates the need for higher-level abstractions that can manage concurrency safely and
declaratively.

\paragraph{The \GRust Kalman filter}
An extract of the \GRust implementation of the Kalman filter, adapted from the \Ci firmware, is
presented in \Cref{lst:grust:benchmarks:compare_grust:kalman}. This version replaces low-level
concurrency primitives with high-level, declarative abstractions.

First, shared memory and manual synchronization are eliminated. Instead, data is exchanged through
the explicit \emph{interface} of a \emph{service}. The complex semaphore mechanism is replaced by
simply exporting the \texttt{estimator\_state} flow. Any downstream service, such as a
\texttt{stabilizer}, can then import this flow and will automatically react to its updates, making
the data dependency explicit and safe.

The core logic of the Kalman filter is modeled as a declarative, event-driven state machine, defined
within the \texttt{when} block (\Crefrange{line:kalman:when:start}{line:kalman:when:end}). This
structure is driven by events derived from input flows. For example, sensor data is modeled as a
\texttt{measure} signal (\Cref{line:kalman:import:meas}). Changes to this signal generate an event
(\Cref{line:kalman:event:meas}) that triggers the main computational step
(\Cref{line:kalman:when:meas}).

Temporal properties are also specified declaratively. The \texttt{@[1, 3000]} annotation
(\Cref{line:kalman:param}) constrains the estimator to react to measurement updates at most once per
millisecond (matching the 1,000~Hz rate of the stabilizer task) and ensures the last known state is
re-emitted if no new data arrives within three seconds.

Within the state machine, specific events trigger corresponding actions. The occurrence of the
\texttt{reset\_ev} event resets the filter's internal state (\Cref{line:kalman:when:reset}), while
an update to the \texttt{measure} signal triggers the prediction and correction sequence
(\Cref{line:kalman:when:meas:predict,line:kalman:when:meas:update}). The model is easily extensible:
additional behaviors can be added simply by defining new events and their corresponding handlers,
without modifying complex, imperative control flow.

\lstinputlisting[
    linerange={1-40},
    language=GRust,
    caption={[Extract of an implementation of a Kalman Filter encoded in \GRust.]%
            Extract of an implementation of the Kalman Filter encoded in \GRust.},
    label={lst:grust:benchmarks:compare_grust:kalman},
    captionpos=b
]{crazyflie_kalman_estimator.rs}

\paragraph{Results}

This case study demonstrates that the \GRust model offers significant improvements in clarity and
conciseness over the original \Ci implementation. By replacing manual synchronization and imperative
control flow with high-level, declarative constructs, the resulting code is not only more compact
but also more readable, as its logic is easier to follow.

However, this example also highlights an area for future improvement in the current compiler
prototype. The syntax of the \texttt{when} equation requires all variables to be declared at the
component level, even if they are only used within a single branch. This leads to the repeated use
of \texttt{let} statements and can result in unnecessary memory allocations, as seen with the
\texttt{predicted} and \texttt{noisy} variables. A more advanced implementation could introduce
block-scoped local variables to address this inefficiency.

Finally, it is important to note that the primary goal of this case study is to evaluate design and
readability, not to compare runtime performance. The original \Ci implementation, with its intricate
use of semaphores, is a testament to the manual optimization required for resource-constrained
microcontrollers. Such low-level code is engineered to extract maximum performance, a goal that
often comes at the expense of clarity and maintainability. A quantitative performance benchmark was
therefore not conducted. The hardware of the \crazyflie is not a suitable target for asynchronous
runtime of \Rust, which is designed for the more powerful computing platforms typical of \sdv
architectures. On these more capable systems, the development and verification cost of manual,
low-level optimization often outweighs the benefits, making a high-level, safety-focused language
like \GRust a more effective choice.

\subsection{Comparing \GRust to Handmade \Rust Code}
\label{sec:grust:benchmarks:grust_vs_rust}

As explained in \Cref{sec:intro:gap}, \ampere is adopting asynchronous \Rust to meet the performance
and concurrency demands of \sdv architectures. However, direct implementation in a low-level
language like \Rust is challenging, and the complexities of asynchronous programming further
increase the risk of errors. \GRust addresses this by automatically generating safe and efficient
asynchronous \Rust code from a high-level model. To validate this approach, we must demonstrate that
the generated code matches the performance of handwritten, manually optimized implementations.
Success would confirm that the abstractions and safety guarantees of \GRust come at no performance
cost, justifying its adoption over complex and error-prone manual coding.

\paragraph{Implemented programs}

For this evaluation, conducted during an internship at \ampere, we assessed the expressiveness,
usability, and performance of \GRust. The intern Matthieu Massif reimplemented three representative
programs from the existing \sdv codebase of \ampere, chosen to cover a range of modeling paradigms%
---event-driven, purely periodic, and mixed behaviors---found in real-world automotive use cases:

\begin{itemize}
    \item The \emph{over-speed alert}, an implementation from which the program in
          \Cref{lst:grust:benchmarks:grust_vs_rust:code_osa} is derived, raises a warning when the
          vehicle speed remains in a specific range for a defined duration. Its behavior is
          well-suited to event-driven modeling, highlighting how \GRust concisely expresses temporal
          conditions with asynchronous triggers.

    \item The \emph{black-ice alert} monitors external temperature to detect potential black-ice
          conditions. This program is implemented as a purely periodic system, where computation
          is triggered at fixed intervals, demonstrating the support in \GRust for traditional,
          time-triggered control logic.

    \item The \emph{trip mileage} tracker accumulates the total travel distance covered during the
          vehicle's lifetime, combining periodic sensor sampling with sporadic storage updates. As a
          mixed system, it requires both periodic sampling and event-driven responses,
          illustrating how \GRust integrates both paradigms.
\end{itemize}

\lstinputlisting[
    linerange={1-50},
    language=GRust,
    caption={[Derived implementation of the Over-Speed Alert in \GRust.]%
            Derived implementation of the Over-Speed Alert program encoded in \GRust.},
    label={lst:grust:benchmarks:grust_vs_rust:code_osa},
    captionpos=b
]{overspeed_alert.rs}

\newcommand{\ratio}[1]{\textcolor{gray}{(#1)}}
\begin{table}[ht!]
    \centering
    \begin{tabular}{lrrr}
        \toprule
        Program      & \ampere & \GRust & Reduction (\%) \\
        \midrule
        Over-speed   & 294     & 86     & 71             \\
        Black-ice    & 165     & 111    & 33             \\
        Trip Mileage & 597     & 240    & 60             \\
        \bottomrule
    \end{tabular}
    \caption[Lines of code for \ampere implementations and their \GRust reimplementations.]%
    {Lines of code (excluding comments) for the original implementations at \ampere and
        their \GRust reimplementations, with the relative reduction in size.}
    \label{tab:grust:benchmarks:grust_vs_rust:loc}
\end{table}

The original programs, implemented in asynchronous \Rust, were reimplemented in \GRust based on
their specifications. We ensured functional equivalence by validating the new implementations
against the existing test suites. As shown in \Cref{tab:grust:benchmarks:grust_vs_rust:loc}, the \GRust versions
consistently reduce code size, with reductions ranging from 33\% to 71\%. The event-driven
\emph{over-speed alert} shows the most significant improvement, requiring less than a third of the
original code, while the \emph{trip mileage} example is also reduced by 60\%.

Notably, the reimplementations were based on the original specifications, which reflect a
periodic mindset rooted in legacy development practices. These specifications do not fully
exploit the event-driven capabilities that \sdv architectures and \GRust support. A key issue is
that such specifications often dictate implementation details rather than defining behavioral
requirements. For instance, the specification for the \emph{black-ice alert} explicitly requires an
implementation based on periodic temperature monitoring. This approach prescribes
\emph{how} the system should work, instead of stating \emph{what} it must achieve. A more
declarative requirement would be, for example, to filter the temperature signal to prevent changes
of more than one degree per second. Prescriptive styles hinder the adoption of more efficient,
event-driven models.

\paragraph{Benchmarks strategy}

To provide an accurate performance comparison, Matthieu Massif, assisted by \ampere engineers,
deployed both the \GRust-generated code and the handwritten \Rust versions on the target hardware:
the computer board inside the vehicle.
Benchmarking on the target is critical because performance on a developer machine rarely translates
to a resource-constrained embedded environment. The target hardware may have a different CPU
architecture, memory hierarchy, and operating system scheduler, all of which can expose bottlenecks
not visible during desktop simulation. This step ensures that our evaluation is grounded in
real-world conditions.

To evaluate performance under realistic conditions, we generated inputs from a 15~minute driving
scenario. This scenario simulates various driving phases---urban, rural, highway, and parking---to
cover a wide range of operational conditions. The three programs under test are part of a larger
\emph{Display Compute} service named \dispute, which is responsible for computing information for
the cockpit display. This service receives the following inputs:
%
\begin{itemize}
    \item the vehicle speed (updated every 20\,ms),
    \item the raw external temperature (updated every 1000\,ms), and
    \item the total accumulated distance (updated every 300\,ms).
\end{itemize}
%
In response, it computes several outputs, including:
%
the over-speed alert,
the black-ice alert,
the filtered external temperature,
the mileage, and
other outputs not activated by this particular trace.
%
From this scenario, we generated two distinct input traces to simulate different bus communication
strategies. The first, a \emph{periodic} trace, mimics a time-triggered architecture by sending
data at the refresh period of each input, regardless of value changes. The second, an
\emph{on-change} trace, represents an event-driven approach where data is sent only when a value
changes. Using both traces allows us to assess the behavior of the system under different
data-loading conditions.

To quantify performance and resource consumption, we measured several key metrics for both the
original \dispute and its reimplemented version, denoted by \disgrust. Both versions were evaluated
using the periodic and on-change traces.

\begin{description}
    \item[Log Volume] We measured the total number of log messages generated by each version to
          assess the reduction in communication overhead.

    \item[CPU Load] This metric measures the percentage of time the CPU is busy. In a real-time
          system, maintaining a low CPU load is essential to ensure that all tasks meet their
          deadlines and to accommodate future functionalities.

    \item[Memory Consumption (PSS/RSS)] We measured the Resident Set Size (RSS), which is the total
          memory a process occupies in RAM, and the Proportional Set Size (PSS), which provides a
          more accurate representation of the memory overhead of a process by accounting for shared
          libraries. Minimizing memory footprint is crucial on embedded systems, where RAM is a
          scarce and costly resource.

    \item[Page Faults] A page fault occurs when a program accesses a memory page that is not
          currently in RAM. Handling these faults requires the operating system to load data from
          storage, introducing significant and unpredictable latency. For real-time systems,
          minimizing page faults is critical for ensuring deterministic performance and preventing
          missed deadlines.

    \item[Execution Tracing] To gain deeper insight into the execution flow, we captured function
          call traces using the \texttt{tracing} crate~\cite{tracing_rust}, a framework for
          instrumenting \Rust programs to emit structured, event-based data. We then analyzed these
          traces with \perfetto~\cite{perfetto}, an open-source tool for performance trace
          visualization. This analysis was particularly insightful for comparing the execution of
          \disgrust under the periodic and on-change traces. It allowed us to quantify how many
          times internal components were executed, thereby directly measuring the computational
          savings from the event-driven approach, which avoids redundant work when input values
          remain unchanged.

    \item[Functional Correctness] We verified functional correctness by comparing the input and
          output logs, ensuring that for a given input, the generated code produced a similar output
          as the original implementation within the same timeframe.
\end{description}

\paragraph{Results}

We now present the results of our performance evaluation, comparing the original \dispute
implementation with our \GRust-based version, \disgrust. The evaluation was conducted using both
periodic and on-change input traces to assess performance under different communication
strategies. The results for each metric are presented in turn.

\subparagraph{Log Volume}

The number of log messages for each flow is detailed in
\Cref{tab:grust:benchmarks:grust_vs_rust:logs}. The most significant finding is the drastic
reduction in total log volume when using the on-change trace. For both \dispute and \disgrust, the
total number of messages is reduced by approximately 50\%, dropping from around 58,000 to 28,000.
This reduction is primarily due to the \emph{vehicle speed} input, which is transmitted nearly three
times less frequently in the on-change scenario, indicating that its value often remains constant
between updates.

\begin{table}[H]
    \centering
    \begin{tabular}{|r|cccc|}
        \hline
        \multirow{2}{*}{Input/output flows} & \multicolumn{2}{c}{periodic} & \multicolumn{2}{c|}{on-change}                        \\\cline{2-5}
                                            & \dispute                     & \disgrust                      & \dispute & \disgrust \\\hline\hline
        Vehicle speed                       & 44405                        & 44344                          & 15315    & 15161     \\
        Raw external temperature            & 891                          & 885                            & 896      & 886       \\
        Total accumulated distance          & 2550                         & 2550                           & 1828     & 1813      \\
        Over-speed alert                    & 1                            & 187                            & 1        & 185       \\
        Black-ice alert                     & 914                          & 963                            & 899      & 966       \\
        Filtered external temperature       & 920                          & 236                            & 898      & 237       \\
        Mileage                             & 3941                         & 3988                           & 3920     & 3990      \\
        Others                              & 4649                         & 4966                           & 4525     & 4731      \\\hline
        Total                               & 58271                        & 58119                          & 28282    & 27969     \\\hline
    \end{tabular}
    \caption[Number of logs for each execution.]%
    {Number of logs for each execution (lasting 15\,min) and each flow.}
    \label{tab:grust:benchmarks:grust_vs_rust:logs}
\end{table}

The results also reveal behavioral differences between the two implementations. For the
\emph{filtered external temperature} output, the \dispute implementation generates a new value for
nearly every input ($\sim$900 logs). In contrast, \disgrust emits a value only when the filtered
temperature actually changes, resulting in a 75\% reduction in logs for this flow (from $\sim$900 to
$\sim$230).
%
Conversely, the \emph{over-speed alert} output shows a significant divergence. The original \dispute
implementation emits the alert only once. The \disgrust version, however, produces around 185 logs.
This difference stems from how liveness properties are handled. In the original system, the
communication middleware is responsible for ensuring the alert persists, a behavior external to the
\dispute model. In \GRust, this temporal constraint is made explicit in the model itself. The
\disgrust service was specified to re-emit the alert after a timeout to ensure liveness, thus making
the intended behavior part of the verifiable model.

For other flows, such as the \emph{black-ice alert} and \emph{mileage}, the log counts are
comparable across both versions and traces. Overall, while the total log volumes are similar, the
\disgrust version demonstrates a more semantically refined, on-change behavior.

\subparagraph{CPU Load}

The CPU load measurements, presented in \Cref{tab:grust:benchmarks:grust_vs_rust:cpu_load_total},
highlight the primary benefit of an event-driven architecture. For both implementations, switching
from the periodic to the on-change trace cuts CPU utilization by more than half. This improvement
is a direct consequence of the reduced data volume shown previously. By processing inputs only when
their values change, the system avoids not only redundant computation but also significant
communication protocol overhead.

\begin{table}[H]
    \centering
    \begin{tabular}{|r|cc|}
        \hline    & periodic & on-change \\\hline
        \dispute  & 1.88     & 0.87      \\
        \disgrust & 2.13     & 0.96      \\\hline
    \end{tabular}
    \caption[Total CPU load of the \dispute and \disgrust processes.]%
    {Total CPU load (\%) of the \dispute and \disgrust processes for each execution.}
    \label{tab:grust:benchmarks:grust_vs_rust:cpu_load_total}
\end{table}

The results also reveal a slight overhead in the \disgrust implementation compared to the original,
approximately 0.25\% in the periodic case and under 0.1\% in the on-change scenario. This
overhead should not be seen as an inherent cost of the \GRust model. The \GRust project was
conducted in parallel with the development of a new \Rust-based architecture at \ampere.
Consequently, the interface between \disgrust and the rest of the system was implemented manually
and may lack production-level optimizations. This hand-crafted integration likely contributes to the
measured overhead. Future work will focus on creating a seamless interface with the \ampere system,
which is expected to reduce these integration costs.
%
Another portion of the overhead can be attributed to the explicit liveness logic for the
\emph{over-speed alert}, which makes temporal requirements part of the verifiable model. The
narrowing of the performance gap in the on-change scenario also suggests that the asynchronous
framework operates more efficiently on the event-driven workloads for which it is designed.

\subparagraph{Memory Consumption (PSS/RSS)}

The memory consumption of the \dispute and \disgrust processes is detailed in
\Cref{tab:grust:benchmarks:grust_vs_rust:pss,tab:grust:benchmarks:grust_vs_rust:rss}. We evaluate
memory usage using two standard metrics: the \emph{Resident Set Size} (RSS) and the
\emph{Proportional Set Size} (PSS). For both processes, the memory footprint remains nearly constant
across the periodic and on-change executions. This confirms that memory consumption is primarily
a static cost determined by the program structure, rather than a dynamic cost influenced by the
workload.

\begin{table}[H]
    \centering
    \begin{tabular}{|r|cc|}
        \hline    & periodic & on-change \\\hline
        \dispute  & 30450    & 30693     \\
        \disgrust & 30781    & 30827     \\\hline
    \end{tabular}
    \caption[PSS of the \dispute and \disgrust processes.]%
    {PSS (KiB) of the \dispute and \disgrust processes for each execution.}
    \label{tab:grust:benchmarks:grust_vs_rust:pss}
\end{table}

\begin{table}[H]
    \centering
    \begin{tabular}{|r|cc|}
        \hline    & periodic & on-change \\\hline
        \dispute  & 34961    & 34751     \\
        \disgrust & 36821    & 36559     \\\hline
    \end{tabular}
    \caption[RSS of the \dispute and \disgrust processes.]%
    {RSS (KiB) of the \dispute and \disgrust processes for each execution.}
    \label{tab:grust:benchmarks:grust_vs_rust:rss}
\end{table}

The RSS metric measures the total physical memory occupied by a process, including shared libraries.
While it provides an upper bound on the memory footprint, it can overstate the cost of a single
process by attributing the full size of shared libraries. The results show that the RSS of \disgrust
is approximately 1800\,KiB (about 5\%) larger than the RSS of \dispute. This difference suggests
that the \disgrust implementation and its dependencies require a larger set of pages to be loaded
into memory.

For a more precise assessment of the overhead introduced by \disgrust, the PSS metric is more
suitable. PSS accounts for shared memory by dividing the size of each shared library by the number
of processes using it. The PSS results reveal a much smaller overhead, ranging from 134 to 331\,KiB.
This increase corresponds to a marginal overhead of 1\% or less compared to the memory usage of the
handwritten \dispute implementation.

\subparagraph{Page Faults}

The page fault data, presented in \Cref{tab:grust:benchmarks:grust_vs_rust:page_faults}, provides
insight into the memory access patterns of both implementations. A page fault occurs when a program
accesses a memory page that is not in its working set. \emph{Major} faults are costly as they
require loading data from disk, while \emph{minor} faults are less expensive but still introduce
latency as the kernel maps an already-resident page.

The results show that major page faults are rare and identical for both implementations, occurring
only a handful of times. This suggests that both programs load their necessary data into memory
efficiently at startup, with no significant differences in disk I/O behavior during execution.
%
A notable difference, however, appears in the number of minor page faults accumulated over the
15~minute execution. \dispute recorded over 100 minor faults in both scenarios, whereas \disgrust
generated only 14 at most. This improvement could stems from the memory allocation strategy of the
generated code. The \disgrust implementation operates exclusively on the stack, which is allocated
once at startup. This design avoids the dynamic memory requests that are a common source of minor
page faults, leading to a more stable and predictable memory access pattern.

\begin{table}[H]
    \centering
    \begin{tabular}{|rr|cc|}
        \hline                     &       & periodic & on-change \\\hline
        \multirow{2}{*}{\dispute}  & Major & 0        & 4         \\
                                   & Minor & 118      & 109       \\
        \multirow{2}{*}{\disgrust} & Major & 0        & 4         \\
                                   & Minor & 10       & 14        \\\hline
    \end{tabular}
    \caption[Major and minor page faults of the \dispute and \disgrust processes]%
    {Number of major and minor page faults of the \dispute and \disgrust processes
        for each execution (lasting 15\,min).}
    \label{tab:grust:benchmarks:grust_vs_rust:page_faults}
\end{table}

While the rate of faults for \dispute is low---averaging less than one every few seconds---the
near-elimination of these events in \disgrust is a notable advantage. For high-assurance real-time
systems, minimizing OS-level interruptions, even minor ones, is a desirable property for ensuring
deterministic performance.

\subparagraph{Tracing with Perfetto}

To move beyond high-level metrics and understand the underlying execution dynamics, we instrumented
the code using the \texttt{tracing} crate~\cite{tracing_rust} and captured detailed execution
traces, which we then visualized with \perfetto~\cite{perfetto}. This analysis provides a granular
view of the execution flow, allowing us to pinpoint the sources of performance differences between
the periodic and on-change traces in the \disgrust process.

\begin{figure}[H]
    \centering
    \includegraphics[width=\textwidth]{benchs/perfetto_explain_trace.png}
    \caption[Detailed view of a \perfetto trace.]%
    {A detailed view of a \perfetto trace for \disgrust, illustrating a typical
        input-processing cycle and asynchronous runtime behavior. The trace displays the
        state (\eg \emph{Running} or \emph{Runnable}) and active job for each of four
        threads. (Function names have been edited for proprietary reasons).}
    \label{fig:grust:benchmarks:grust_vs_rust:perfetto:zoomed}
\end{figure}

\Cref{fig:grust:benchmarks:grust_vs_rust:perfetto:zoomed} shows a detailed view of a \perfetto
trace. Each horizontal track represents a thread, displaying its state over time---such as
\emph{Runnable} (light green), when ready to execute, or \emph{Running} (dark green), when actively
executing a job. Jobs, which correspond to instrumented function calls, are shown as stacked
rectangles, with the length of the top-level rectangle indicating the total duration of the task.

The second thread illustrates a complete input-processing cycle: a brief task handles the reception
of an input (lilac and dark blue), which triggers a computational step (purple), and finally, a
series of tasks publish the outputs to the communication bus (petrol blue). Visually, the duration
of the publishing tasks is significantly longer than that of the core computational step, confirming
that bus communication is a major contributor to execution time. The trace also reveals dynamics of
the asynchronous runtime, such as one thread being idle (thread four) while another executes work
for a different process (thread three), a behavior typical of work-stealing schedulers.

\begin{figure}[H]
    \centering
    \begin{subfigure}{0.8\linewidth}
        \includegraphics[width=\textwidth]{benchs/perfetto_sync_full_trace.png}
        \caption[Full \perfetto trace for the periodic trace of \disgrust.]%
        {Full \perfetto trace for the periodic trace of \disgrust.}
        \label{fig:grust:benchmarks:grust_vs_rust:perfetto:full:sync}
    \end{subfigure}
    \hfill
    \begin{subfigure}{0.8\linewidth}
        \includegraphics[width=\textwidth]{benchs/perfetto_onchange_full_trace.png}
        \caption[Full \perfetto trace for the on-change trace of \disgrust.]%
        {Full \perfetto trace for the on-change trace of \disgrust.}
        \label{fig:grust:benchmarks:grust_vs_rust:perfetto:full:onchange}
    \end{subfigure}
    \caption[Full \perfetto traces of the \disgrust process.]%
    {Full \perfetto traces of the \disgrust process, comparing the total workload under the
        periodic and on-change scenarios. The dense activity in the periodic trace contrasts
        with the sparser execution of the on-change trace.}
    \label{fig:grust:benchmarks:grust_vs_rust:perfetto:full}
\end{figure}

A high-level comparison of the full execution traces, shown in
\Cref{fig:grust:benchmarks:grust_vs_rust:perfetto:full}, offers immediate visual evidence of the
benefits of an event-driven approach. The trace for the periodic scenario
(\Cref{fig:grust:benchmarks:grust_vs_rust:perfetto:full:sync}) shows dense activity, whereas the
on-change trace (\Cref{fig:grust:benchmarks:grust_vs_rust:perfetto:full:onchange}) is much sparser.
This confirms that the on-change strategy imposes a considerably lower computational demand on the
system.

\begin{figure}[H]
    \centering
    \begin{subfigure}{0.8\linewidth}
        \includegraphics[width=\textwidth]{benchs/perfetto_sync_slice_flamegraph.png}
        \caption[\perfetto \emph{Slice FlameGraph} for the periodic trace.]
        {\perfetto \emph{Slice FlameGraph} for the periodic trace of \disgrust.}
        \label{fig:grust:benchmarks:grust_vs_rust:perfetto:slice:sync}
    \end{subfigure}
    \hfill
    \begin{subfigure}{0.8\linewidth}
        \includegraphics[width=\textwidth]{benchs/perfetto_onchange_slice_flamegraph.png}
        \caption[\perfetto \emph{Slice FlameGraph} for the on-change trace.]%
        {\perfetto \emph{Slice FlameGraph} for the on-change trace of \disgrust.}
        \label{fig:grust:benchmarks:grust_vs_rust:perfetto:slice:onchange}
    \end{subfigure}
    \caption[\perfetto \emph{Slice FlameGraph} visualizations.]%
    {\perfetto \emph{Slice FlameGraph} visualizations for the periodic and on-change
        traces of \disgrust. These graphs aggregate function calls, with block width proportional to
        the call count, to profile the most frequent operations. (Function names have been edited
        for proprietary reasons).}
    \label{fig:grust:benchmarks:grust_vs_rust:perfetto:slice}
\end{figure}

\begin{table}[H]
    \centering
    \begin{tabular}{|r|cc|}
        \hline
        Slice name                   & periodic samples                & on-change samples              \\\hline
        \texttt{root}                & 12724 \textcolor{gray}{(100\%)} & 8771 \textcolor{gray}{(100\%)} \\
        \texttt{do\_processing}      & 2128 \textcolor{gray}{(17\%)}   & 1125 \textcolor{gray}{(13\%)}  \\
        \texttt{observe}             & 4216 \textcolor{gray}{(33\%)}   & 2240 \textcolor{gray}{(25\%)}  \\
        \texttt{subscriber variants} & 2111 \textcolor{gray}{(17\%)}   & 1121 \textcolor{gray}{(13\%)}  \\
        \texttt{steps}               & 1294 \textcolor{gray}{(10\%)}   & 1324 \textcolor{gray}{(15\%)}  \\
        \texttt{publisher variants}  & 425 \textcolor{gray}{(3\%)}     & 423 \textcolor{gray}{(5\%)}    \\
        \texttt{publish}             & 2550 \textcolor{gray}{(20\%)}   & 2538 \textcolor{gray}{(29\%)}  \\\hline
    \end{tabular}
    \caption[Function call counts, from the \perfetto \emph{Slice FlameGraph}.]%
    {Function call counts extracted from the \perfetto \emph{Slice FlameGraph} for both
        traces. The table quantifies the reduction in input-handling operations (\eg
        \texttt{observe}) \versus the stable number of core computations (\texttt{steps}) in the
        on-change scenario.}
    \label{tab:grust:benchmarks:grust_vs_rust:perfetto:samples}
\end{table}

To quantify these differences, we use the \emph{Slice FlameGraph} view in \perfetto, shown in
\Cref{fig:grust:benchmarks:grust_vs_rust:perfetto:slice}. This visualization aggregates all timed
sections (\textit{slices}) from the trace and groups them by name. The width of each block in the
resulting graph is proportional to the total number of times that function was called, providing a
clear profile of the most frequent operations.

The data extracted from these flame graphs, presented in
\Cref{tab:grust:benchmarks:grust_vs_rust:perfetto:samples}, confirms our hypothesis. The number of
calls to input-handling functions (\texttt{do\_processing}, \texttt{observe}, and \texttt{subscriber
    variants}) is nearly halved in the on-change trace, a direct reflection of the reduction in
redundant input arrivals. In contrast, the number of calls to the core computational logic
(\texttt{steps}) and output-publishing functions (\texttt{publish}) remains almost identical. This
result highlights the efficiency of the \GRust execution strategy: computations are performed only
when an input to a step function changes, and outputs are sent only when their values change.

Notably, the relative contribution of \texttt{publish} calls rises from 20\% to 29\% of the total
samples in the on-change scenario. This does not mean publishing became more costly; rather, as the
total number of input-related calls decreased, the constant number of publishing calls became a more
dominant part of the remaining work. This analysis makes it clear that the on-change strategy saves
CPU time by eliminating the overhead of processing unchanged inputs. Since the core logic is
executed only when necessary, the threads have more idle time, leading to the significant reduction
in overall CPU load.

\subparagraph{Functional Correctness}

To verify functional equivalence, we compared the output values generated by \dispute and \disgrust
for identical input streams. The statistical summary of the relative differences is presented in
\Cref{tab:grust:benchmarks:grust_vs_rust:diff:onchange} for the on-change trace and
\Cref{tab:grust:benchmarks:grust_vs_rust:diff:periodic} for the periodic trace.

The results demonstrate a high degree of functional equivalence. For most outputs, the mean and
median differences are zero, indicating that both implementations produce identical values for the
vast majority of the execution. The small, non-zero discrepancies observed can be attributed to
minor timing variations in the on-board logging mechanism.

Two notable anomalies warrant further explanation. The first is the consistent, non-zero difference
in the \emph{filtered external temperature} and, consequently, the \emph{black-ice alert}. This
divergence is not a functional bug but a direct result of different initial state values for the
temperature filter. As illustrated in \Cref{fig:grust:benchmarks:grust_vs_rust:compare:temp}, the
\dispute implementation starts with a low initial temperature, while \disgrust starts high. Both
traces then converge to the same value after approximately 5~minutes. For the remainder of the
execution, their outputs are nearly identical. This confirms that while their initial states differ,
their dynamic filtering logic is functionally equivalent.
%
The second anomaly is the large maximum error for the \emph{mileage} output in the on-change trace.
This is not a persistent error but an isolated singularity, as shown in
\Cref{fig:grust:benchmarks:grust_vs_rust:compare:mileage:onchange}. Further analysis, presented in
\Cref{fig:grust:benchmarks:grust_vs_rust:compare:mileage:dispute}, reveals that this spike
originates from a singularity within the original \dispute implementation, which manifests
only under the on-change trace. This type of subtle, input-dependent runtime error highlights a
critical challenge in developing high-assurance software. It is precisely this class of defect that
formal verification aims to eliminate. Using a tool like \GReusot, a developer could formally prove
at the model level the absence of such singularities.

Therefore, after accounting for initialization differences, minor logging artifacts, and the
identified singularity, we conclude that the core logic of the \disgrust implementation is
functionally equivalent to the original \dispute service for this scenario.

\begin{figure}[H]
    \centering
    \begin{subtable}{\linewidth}
        \centering
        \footnotesize
        \begin{tabular}{|r|c|c|ccc|c|}
            \hline
            \multirow{2}{*}{Input/output flows}         & \multirow{2}{*}{Mean} & \multirow{2}{*}{Median} & \multicolumn{3}{c|}{Quantiles} & \multirow{2}{*}{Max}                                            \\\cline{4-6}
                                                        &                       &                         & \textonequarter                & \multicolumn{1}{|c|}{\textonehalf} & \textthreequarters &       \\\hline\hline
            Vehicle speed                 ($\pm 0.001$) & 0.005                 & 0.004                   & 0.000                          & \multicolumn{1}{|c|}{0.004}        & 0.007              & 2.000 \\
            Raw external temperature      ($\pm 0.001$) & 0.000                 & 0.000                   & 0.000                          & \multicolumn{1}{|c|}{0.000}        & 0.000              & 0.005 \\
            Total accumulated distance    ($\pm 0.001$) & 0.000                 & 0.000                   & 0.000                          & \multicolumn{1}{|c|}{0.000}        & 0.001              & 0.002 \\
            Over-speed alert              ($\pm 0.001$) & 0.000                 & 0.000                   & 0.000                          & \multicolumn{1}{|c|}{0.000}        & 0.000              & 0.000 \\
            Black-ice alert               ($\pm 0.001$) & 0.062                 & 0.000                   & 0.000                          & \multicolumn{1}{|c|}{0.000}        & 0.000              & 0.667 \\
            Filtered external temperature ($\pm 0.001$) & 0.203                 & 0.000                   & 0.000                          & \multicolumn{1}{|c|}{0.000}        & 0.016              & 1.529 \\
            Mileage                       ($\pm 0.001$) & 0.007                 & 0.000                   & 0.000                          & \multicolumn{1}{|c|}{0.000}        & 0.000              & 2.000 \\\hline
        \end{tabular}
        \caption[Relative difference between \dispute and \disgrust with the on-change trace.]%
        {Summary of the relative difference (\%) between \dispute and \disgrust using the
            on-change trace.}
        \label{tab:grust:benchmarks:grust_vs_rust:diff:onchange}
    \end{subtable}
    \hfill
    \begin{subtable}{\linewidth}
        \centering
        \footnotesize
        \begin{tabular}{|r|c|c|ccc|c|}
            \hline
            \multirow{2}{*}{Input/output flows}         & \multirow{2}{*}{Mean} & \multirow{2}{*}{Median} & \multicolumn{3}{c|}{Quantiles} & \multirow{2}{*}{Max}                                            \\\cline{4-6}
                                                        &                       &                         & \textonequarter                & \multicolumn{1}{|c|}{\textonehalf} & \textthreequarters &       \\\hline\hline
            Vehicle speed                 ($\pm 0.001$) & 0.006                 & 0.000                   & 0.000                          & \multicolumn{1}{|c|}{0.000}        & 0.007              & 2.000 \\
            Raw external temperature      ($\pm 0.001$) & 0.000                 & 0.000                   & 0.000                          & \multicolumn{1}{|c|}{0.000}        & 0.000              & 0.005 \\
            Total accumulated distance    ($\pm 0.001$) & 0.000                 & 0.000                   & 0.000                          & \multicolumn{1}{|c|}{0.000}        & 0.000              & 0.002 \\
            Over-speed alert              ($\pm 0.001$) & 0.000                 & 0.000                   & 0.000                          & \multicolumn{1}{|c|}{0.000}        & 0.000              & 0.000 \\
            Black-ice alert               ($\pm 0.001$) & 0.063                 & 0.000                   & 0.000                          & \multicolumn{1}{|c|}{0.000}        & 0.000              & 0.667 \\
            Filtered external temperature ($\pm 0.001$) & 0.195                 & 0.000                   & 0.000                          & \multicolumn{1}{|c|}{0.000}        & 0.000              & 1.529 \\
            Mileage                       ($\pm 0.001$) & 0.000                 & 0.000                   & 0.000                          & \multicolumn{1}{|c|}{0.000}        & 0.000              & 0.001 \\\hline
        \end{tabular}
        \caption[Relative difference between \dispute and \disgrust with the periodic trace.]%
        {Summary of the relative difference (\%) between \dispute and \disgrust using the
            periodic trace.}
        \label{tab:grust:benchmarks:grust_vs_rust:diff:periodic}
    \end{subtable}
    \begin{subfigure}[t]{0.48\textwidth}
        \centering
        \includegraphics[width=\linewidth]{benchs/compare_filtered_temp.png}
        \caption[Filtered temperature of \dispute and \disgrust using the on-change trace.]%
        {Filtered temperature of \dispute and \disgrust using the on-change trace.}
        \label{fig:grust:benchmarks:grust_vs_rust:compare:temp}
    \end{subfigure}
    \\
    \begin{subfigure}[t]{0.48\textwidth}
        \centering
        \includegraphics[width=\linewidth]{benchs/compare_mileage_onchange.png}
        \caption[Mileage of \dispute and \disgrust using the on-change trace.]%
        {Mileage of \dispute and \disgrust using the on-change trace.}
        \label{fig:grust:benchmarks:grust_vs_rust:compare:mileage:onchange}
    \end{subfigure}
    \hfill
    \begin{subfigure}[t]{0.48\textwidth}
        \centering
        \includegraphics[width=\linewidth]{benchs/compare_mileage_dispute.png}
        \caption[Mileage of \dispute using periodic and on-change traces.]%
        {Mileage of \dispute using periodic and on-change traces.}
        \label{fig:grust:benchmarks:grust_vs_rust:compare:mileage:dispute}
    \end{subfigure}
    \caption[Functional comparison of \dispute and \disgrust.]%
    {Comparison of \dispute and \disgrust using periodic and on-change traces.}
\end{figure}

\clearpage

\subparagraph{Experimental Validation of the Oversampling Theorem}

We provide an experimental validation of the oversampling
\Cref{thm:semantics:grfrp:theorems:oversampling}, formally established in
\Cref{sec:semantics:grfrp:theorems:oversampling}. This theorem provides the theoretical foundation
for a core benefit of \GRust: enabling efficient, event-driven execution without sacrificing
correctness. It guarantees that the functional behavior of a service is not altered by redundant
computations, which occur when inputs are processed despite their values not having changed. To
verify this property in practice, we compared the output traces of the \disgrust implementation when
executed with the on-change and the periodic (oversampling) input streams.

\begin{table}[H]
    \centering
    \begin{tabular}{|r|c|c|ccc|c|}
        \hline
        \multirow{2}{*}{Input/output flows}         & \multirow{2}{*}{Mean} & \multirow{2}{*}{Median} & \multicolumn{3}{c|}{Quantiles} & \multirow{2}{*}{Max}                                            \\\cline{4-6}
                                                    &                       &                         & \textonequarter                & \multicolumn{1}{|c|}{\textonehalf} & \textthreequarters &       \\\hline\hline
        Vehicle speed                 ($\pm 0.001$) & 0.037                 & 0.006                   & 0.000                          & \multicolumn{1}{|c|}{0.006}        & 0.016              & 2.000 \\
        Raw external temperature      ($\pm 0.001$) & 0.001                 & 0.000                   & 0.000                          & \multicolumn{1}{|c|}{0.000}        & 0.000              & 0.005 \\
        Total accumulated distance    ($\pm 0.001$) & 0.005                 & 0.005                   & 0.004                          & \multicolumn{1}{|c|}{0.005}        & 0.006              & 0.009 \\
        Over-speed alert              ($\pm 0.001$) & 0.000                 & 0.000                   & 0.000                          & \multicolumn{1}{|c|}{0.000}        & 0.000              & 0.000 \\
        Black-ice alert               ($\pm 0.001$) & 0.000                 & 0.000                   & 0.000                          & \multicolumn{1}{|c|}{0.000}        & 0.000              & 0.000 \\
        Filtered external temperature ($\pm 0.001$) & 0.005                 & 0.000                   & 0.000                          & \multicolumn{1}{|c|}{0.000}        & 0.000              & 0.182 \\
        Mileage                       ($\pm 0.001$) & 0.008                 & 0.005                   & 0.004                          & \multicolumn{1}{|c|}{0.005}        & 0.006              & 2.000 \\\hline
    \end{tabular}
    \caption[Relative difference between the periodic and on-change traces for \disgrust.]%
    {Summary of the relative output difference (\%) between the periodic and on-change
        traces for \disgrust.}
    \label{tab:grust:benchmarks:grust_vs_rust:diff:disgrust}
\end{table}

The results, summarized in \Cref{tab:grust:benchmarks:grust_vs_rust:diff:disgrust}, strongly support
the theorem. For all outputs, the mean and median differences are effectively zero, confirming that
the efficient, event-driven execution produces the same functional result as the traditional
periodic one.

\begin{figure}[H]
    \centering
    \includegraphics[width=0.8\linewidth]{benchs/compare_mileage_disgrust.png}
    \caption[Comparison of the mileage output for the periodic and on-change traces.]%
    {Comparison of the mileage output from \disgrust for the periodic and on-change
        traces, showing functionally identical behavior with minimal latency differences at the
        initial jump.}
    \label{fig:grust:benchmarks:grust_vs_rust:compare:mileage:disgrust}
\end{figure}

The large maximum difference recorded for the \emph{mileage} output does not contradict this
conclusion. This statistical anomaly is a predictable artifact of comparing the two execution
models, and is distinct from the numerical instability found in the original \dispute
implementation. As illustrated in
\Cref{fig:grust:benchmarks:grust_vs_rust:compare:mileage:disgrust}, the two traces are functionally
identical. The momentary divergence occurs because the precise timings of computations are
different. In the brief interval between these two computations, the time-aligned comparison
registers a transient spike, which is captured as the maximum value.

This experiment thus provides strong evidence that the semantics of our model are correctly
implemented, ensuring that the performance gains from on-change execution do not come at the cost of
functional correctness.

\paragraph{Analysis}

The comprehensive evaluation demonstrates that \GRust successfully generates code matching the
performance of handwritten, manually optimized \Rust, while offering significant advantages in
conciseness, maintainability, and correctness.

Performance metrics confirm that the abstractions provided by \GRust introduce no significant
overhead. The slight increase in CPU load and memory footprint observed in the \disgrust
implementation could be a result of the manually implemented interface and the explicit liveness
properties---a design choice that integrates temporal requirements into the verifiable model.
Furthermore, the evaluation revealed a notable improvement in memory behavior. The near-elimination
of minor page faults indicates a more predictable, stack-based memory allocation strategy, which is
highly desirable for ensuring deterministic execution in real-time systems.

The benchmark also highlights the benefits of the on-change execution strategy, which reduced the
CPU load by half for both implementations. \perfetto traces confirmed that this reduction
stems directly from minimizing communication protocol activity, which was identified as the primary
source of CPU consumption. By eliminating redundant data transfers, this approach effectively
targets the main performance bottleneck. Crucially, the experimental validation of the oversampling
theorem provides assurance that these performance gains do not compromise functional correctness.
This result validates the event-driven approach for \sdv architectures and positions \GRust as an
effective tool for its implementation.

In conclusion, the results confirm that \GRust successfully bridges the gap between high-level
modeling and the \Rust ecosystem. It delivers the performance of low-level, asynchronous code while
abstracting its complexity. This offers a clear path toward improved productivity, maintainability,
and safety in software development for \sdv.

\subsection{Summary}
\label{sec:grust:benchmarks:summary}

This section presented a three-part evaluation to validate the core contributions of \GRust: the
performance of its synchronous kernel, its expressiveness for high-level design, and its practical
applicability in an industrial \sdv context.

First, a quantitative benchmark of the synchronous kernel, \GRSync, was performed on systems ranging
from small academic programs to realistic embedded controllers and a semi-industrial aircraft model.
This evaluation showed that the \Rust-based toolchain is a highly competitive solution for
implementing synchronous systems.

Second, a qualitative case study contrasted a \GRust model of a UAV's Kalman filter with its
original, low-level \Ci implementation. The comparison demonstrated that \GRust replaces complex and
error-prone manual concurrency management with high-level, declarative abstractions. This shift
improves design clarity, simplifies the expression of temporal and data dependencies, and enhances
maintainability and safety.

Finally, an industrial evaluation at \ampere compared \GRust-generated code against handwritten,
asynchronous \Rust for three real-world automotive services. The \GRust models were significantly
more concise, reducing code size by up to 71\%. On target hardware, the generated code matched the
performance of the manually optimized versions, with negligible overhead and more predictable memory
behavior. This evaluation also confirmed the benefits of an event-driven approach, which halved CPU
load, and experimentally validated our oversampling
\Cref{thm:semantics:grfrp:theorems:oversampling}, ensuring that these performance gains do not
compromise functional correctness.

Together, these evaluations demonstrate that \GRust successfully bridges the gap between high-level
modeling and low-level implementation. It offers a clear path to developing complex automotive
software with improved productivity and safety, without sacrificing the performance required for
modern \sdv architectures.
